\subsection{Apple: debt balances and cash repayment}

Apple's borrowing shows why the distinction between a contract and its carrying
amount matters for a funding forecast. A fall in reported debt might suggest
how much cash the company repaid. Its annual disclosure lets us test that
interpretation by separating principal from the adjustments included in the
carrying amount. The following amounts are in millions of US dollars at the
displayed precision; changes run from 30 September 2023 to 28 September 2024
\citep[Note 9, printed pp.~43--44]{apple2024}.

\begin{center}
\begin{tabular}{lrrr}
\toprule
Term-debt component & 2023 & 2024 & Change\\\midrule
Principal &106,572&97,341&$-9,231$\\
Net premium, discount and issuance costs &$-356$&$-321$&$+35$\\
Hedge accounting fair-value adjustment &$-1,113$&$-358$&$+755$\\\midrule
Carrying amount &105,103&96,662&$-8,441$\\\bottomrule
\end{tabular}
\end{center}

The cash-flow statement reports no net term-debt issuance proceeds and cash
repayments of 9,958 million for fiscal 2024 \citep[printed p.~33]{apple2024}.
To connect that cash movement to the table, first calculate the difference
between the change in reported principal and net term-debt cash flow. It is
$-9{,}231-(-9{,}958)=727$ million. Adding the reported measurement-component
changes then gives the carrying-amount movement. With $\Delta$ denoting closing
less opening balance, the two displayed-amount identities are
\[
\begin{aligned}
 \Delta\text{principal}
   &=-9{,}958+727=-9{,}231,\\
 \Delta\text{carrying amount}
   &=-9{,}958+727+35+755=-8{,}441.
\end{aligned}
\]
The 727 is an unallocated balance-to-cash difference. It does not identify an
exchange movement or a journal. Likewise, the 35 is the change in a net
unamortised balance, rather than an identified issuance fee or interest charge;
the 755 is a change in the debt's hedge adjustment, rather than a cash receipt.
Together they explain 790 of the 1,517 million difference between the
carrying-amount movement and net cash repayment. A forecast that used the
8,441 carrying decline as repayment would therefore compute the wrong cash
amount for this observed year. The cash statement's proceeds measure is net;
it must not automatically become legal face issued in another period.

The maturity schedule answers a different funding question. Its six future
principal buckets sum to 97,341 million, with 10,930 in fiscal 2025. The current
carrying amount is 10,912. Their 18 million difference does not allocate a
discount or hedge adjustment to the current portion. The measurements differ,
and their time boundaries must also be retained: Note 1 supplies the fiscal-year
default, the issuance table explicitly labels its 2024--2062 range as calendar
years, and MD\&A separately reports amounts payable within twelve months.
MD\&A also reports future term-debt interest payments of 38.5 billion dollars,
including 2.6 billion within twelve months \citep[printed pp.~25, 34 and 43]{apple2024}.
Those cash requirements cannot replace period interest expense or a present
value. The dated Level 2 fair-value disclosures likewise measure something
different: the annual note reports 88.4 billion dollars at September 2024,
while the interim note reports 82.7 billion at December 2024 alongside a
rounded carrying amount of 94.8 billion
\citep[Note 6, printed p.~10]{apple2025q1}.

\subsection{Net borrowing can conceal refinancing}

Commercial paper exposes a second limitation of aggregate financing data.
Apple describes unsecured short-term promissory notes used for general
corporate purposes, including dividends and repurchases. The annual note says
their maturities are generally less than nine months; that description does
not identify each note or trace individual borrowed dollars to a use.
At the annual statement amounts, commercial paper increases from 5,985 to
9,967 million, a rise of 3,982, while net cash proceeds are 3,960. The resulting
22 million difference remains unallocated. During the following quarter the
balance falls from 9,967 to 1,995, while net cash repayments are 7,944; the
corresponding unallocated difference is $-28$ million
\citep[Note 9 and balance and cash-flow statements]{apple2024}
\citep[Note 6 and balance and cash-flow statements]{apple2025q1}.

Even a complete explanation of those differences would leave refinancing
activity unresolved. Annual Note 9 reports net cash for paper with maturities
of at most 90 days, but separate proceeds and repayments for longer maturities.
In fiscal 2022, for example, the longer-maturity proceeds of 5,948 and
repayments of 7,257 give net repayments of 1,309 million. Adding short-maturity
net proceeds of 5,264 gives total net proceeds of 3,955. For fiscal 2024 the
longer-maturity cash rows are zero and the short-maturity net inflow is 3,960.
Neither that net amount nor a closing balance reveals how many issues and
redemptions occurred between reporting dates.

Consider a hypothetical incremental borrowing of 1,000 million US dollars,
issued at par and repaid after 60 days. Stipulate simple interest of 5\% using
actual days over a 360-day year, with no fees, taxes, hedges or early settlement.
The rate is an assumption of this example; Apple's disclosed 5.00\% weighted
rate at September 2024 does not establish the rate on such a transaction.
Issuing and repaying the principal adds $1{,}000-1{,}000=0$ to net principal
cash flow and leaves zero principal outstanding, yet supplies 1,000 million
between those dates. The cost over the 60 days is
\[
 \text{interest}=1{,}000\times0.05\times\frac{60}{360}
               =\frac{25}{3}\approx8.3333
               \quad\text{million US dollars}.
\]
Issue debits cash and credits principal 1,000. Accrual debits interest expense
and credits interest payable $25/3$. Settlement debits principal 1,000 and
interest payable $25/3$, and credits cash $3{,}025/3$. Each entry balances;
after settlement the incremental liability is zero, while cash and equity are
each lower by $25/3$ relative to the otherwise funded baseline.

This example isolates what a net principal-flow observation fails to reveal.
It is not a reconstructed Apple history: interest expense, interest cash and
other disclosures could distinguish or exclude the added borrowing. Within
that scope, it shows why a financing model needs dated issuance, accrual and
redemption events. Matching a net financing row does not establish either
intraperiod liquidity or funding cost.

\subsection{A rate scenario needs an exposure definition}

The contract's coupon is also insufficient to describe exposure after hedging.
Apple's annual debt note states that interest-rate swaps effectively convert
some fixed-rate US-dollar notes to floating rates, while cross-currency swaps
effectively convert some foreign-currency notes to US-dollar notes. The
effective-rate ranges include interest, discount or premium amortisation and
applicable hedge adjustments. These are relationships between separate
contracts and measurements, not a single rate that can be applied to total
principal \citep[Note 9, printed p.~43]{apple2024}.

The filing supplies a useful, more specific constraint. Annual Item 7A reports
that a hypothetical increase of 100 basis points at all tenors, including
associated derivatives, would increase annual term-debt interest expense by
139 million dollars at September 2024 and 194 million at September 2023
\citep[printed p.~27]{apple2024}. Multiplying Note 4's 2024 interest-rate
contract notional of 14,575 million by 1\% instead gives 145.75 million. It does
not reproduce the reported sensitivity, even allowing for nearest rounding of
the displayed amounts. The difference cannot be assigned to a particular
reset date or swap feature from these aggregates. For 2023 the same shortcut
gives $19{,}375\times0.01=193.75$, consistent with rounded 194; that numerical
match does not validate the shortcut or recover the issuer's model.

The scenario's measure matters as much as its shock. In the same table, the
2024 investment-portfolio response is a 2,755 million decline in fair value.
It cannot be added to an annual interest-expense increase as though both were
the same loss measure. At December 2024 the reported interest-rate notional
is still 14,575 million, while the carrying amount of term debt subject to
fair-value hedges is a rounded 13.3 billion, compared with 13.5 billion at
September. These aggregates describe different populations or measurements.
The interim statement that market risk has not changed materially refers the
reader to annual Item 7A; it does not publish an updated 139 million sensitivity
\citep[Note 4 and Part I, Item 3]{apple2025q1}.

Hedge accounting and cash presentation also remain distinct. Annual Note 1
presents derivative assets and liabilities at gross fair values. Note 4
classifies derivative cash flows with the items hedged, which are
\emph{generally} operating activities; that qualifier does not classify every
derivative settlement as operating. The maximum disclosed horizon for hedging
debt-related foreign-currency cash-flow variability is 18 years, and Apple
states that some exposures may be left unhedged and that offsets may be partial
\citep[Notes 1 and 4, printed pp.~34 and 37--38]{apple2024}.

Finally, Item 7A's currency-risk statistic has its own population. Apple reports
a Monte Carlo estimate, at 95\% confidence over one day, of 538 million dollars
at September 2024 and 669 million at September 2023 for adverse fair-value
changes in its foreign-currency derivatives. The assessment excludes underlying
currency exposures, whose changes generally offset derivative changes. This
derivative VaR is neither a hard maximum loss nor a measure of the company's
net currency risk; the issuer's simulation has not been replicated here.
A financing simulation can retain these dated scenarios as disclosed
constraints, but still needs actual holdings, currencies, payment calendars,
reset terms and hedge relationships to calculate future cash and accounting
consequences. The next question is how financing choices change the quantities
and claims of the company's owners.
